\begin{frame}{``앞쪽 층은 범용하다''는 말의 첫 증거}
  \begin{itemize}
    \item 추측이 아니라 \textbf{보여준 것} --- AlexNet 논문
      (Krizhevsky 등, 2012)이 학습된 \textbf{첫 합성곱 층의
      필터를 그림으로 인쇄}했고, 거기서는 방향/색 대비를
      감지하는 \textbf{Gabor 같은 필터}들이 늘어섬
    \item 장 머리에서 허블·비셀이 고양이 시각피질에서 본
      \textbf{단순 세포의 수용장과 동일한 종류}의 패턴
    \item ImageNet 1,000개 클래스를 구분하기 위해 학습된
      네트워크의 첫 층이 ``고양이 특화''가 아니라
      \textbf{``모서리 감지''}에 쓰였다는 뜻
  \end{itemize}
  \vskip0.5em
  \begin{block}{전이학습이 이용하는 비대칭}
    층이 깊어질수록 특징이 눈·바퀴·사물 부분으로
    ``사물화''되는 계층 구조(10.2) 때문에,
    \textbf{앞으로 갈수록 재사용 가치가 높다} ---
    \textbf{앞쪽 = 범용, 뒤쪽 = 특화}.
  \end{block}
\end{frame}
